\section{背景与相关工作}
\label{sec:related}

本节先对所要求解的测量量作形式化定义，并据此确立筛选相关方法的判据；随后依次讨论
在该判据下真正可比的方法及其局限，最后定位本文所填补的空白。

\subsection{问题定义与筛选判据}

给定空间中一条由激光定义的三维直线（激光基准轴 $\mathcal{L}$）与一个目标点
$P_T$，所要求解的物理量是 $P_T$ 到 $\mathcal{L}$ 的垂直距离 $H$。在大型转动设备
对中场景中，$\mathcal{L}$ 对应替代钢丝的理想基准轴，$P_T$ 对应待对中轴上的采样
点，$H$ 即该点相对基准的偏差。典型的精度需求为 $\pm15\,\mu\mathrm{m}$（理想）
到 $\pm50\,\mu\mathrm{m}$（工程可接受），量程为数米到数十米。

判断一个已有方法是否与本文相关，标准只有一个：它能否（直接或间接）给出“基准轴外
一点到该基准轴的距离”这一几何量。至于方法内部是否用到激光、其激光扮演什么角色，
与本判据无关——我们关心的是测量方法本身能否迁移到本文场景，而非其实现细节。据此，
一大类看似涉及“激光”的仪器其实并不可比：激光三角测距测的是传感器本体到目标表面
沿其光轴的单点位移，激光干涉测长测的是随动光学部件相对光路的位移，自准直仪测的
则是反射镜的角度。它们的输出量都不是“点到某条外部基准轴的距离”，也无法在不引入
额外基准定义的情况下迁移到本场景，因此不在本文的对比范围之内。下面只讨论在上述
判据下确实能测 $H$（或其变体）的方法，如表~\ref{tab:compare} 所示。

\begin{table*}[!t]
\centering
\caption{可测“点到基准轴距离”的方法对比}
\label{tab:compare}
\begin{tabular}{@{}p{2.6cm}p{3.0cm}p{2.4cm}p{4.2cm}@{}}
\toprule
方法 & 代表产品／文献 & 测 $H$ 的方式 & 主要局限 \\
\midrule
钢丝找中法 & 传统工装~\cite{wire_alignment} &
直接 & 钢丝下坠、振动敏感、需反复拆装、量程受限 \\
激光跟踪仪 & API Radian、Leica AT960~\cite{api_radian,leica_at960} &
间接（测坐标再算距离） & 需合作反射器，无法测非合作点，单价高 \\
双目激光线重建 & Zhong~\cite{zhong2021line3d}、Xu／Chen~\cite{xu2017ortho,chen2021vision2cam} &
直接（重建线／点） & 精度仅毫米级，达不到工程要求 \\
PSD／靶板对中仪 & Easy-Laser XT770、Fixturlaser NXA Pro~\cite{easylaser_xt770,fixturlaser_nxa} &
仅限受测单点 & 只能测激光斑击中探测器的那一点，换点须移装 \\
\textbf{DCPAM（本文）} & 双相机 + 接收屏 + 镜面虚像重建 &
\textbf{直接（对任意点）} & 依赖 PnP 与镜面装配标定精度 \\
\bottomrule
\end{tabular}
\end{table*}

\subsection{传统钢丝找中法}

在所有方法中，唯一与 $H$ 语义完全一致、且现场长期沿用的是传统钢丝找中法~\cite{wire_alignment}。
它以一根张紧钢丝作基准直线，逐点测量轴表面到钢丝的垂直距离，测的正是“点到基准轴
距离”，因而是本文最直接的对标对象，也是要替代的方法。其局限已在第~\ref{sec:intro}
节详述：钢丝自重下坠使基准并非真正的直线、对气流与振动敏感、且每轮测量都需反复
架设拆装。DCPAM 保留了“点到基准轴距离”这一测量语义，同时用激光基准消除下坠、
用非接触方式免除反复拆装。

\subsection{激光跟踪仪与全站仪}

激光跟踪仪用角度编码器加测距单元跟踪一个合作反射器（SMR），输出反射器球心的三维
坐标，代表产品如 API Radian~\cite{api_radian} 与 Leica Absolute Tracker
AT960~\cite{leica_at960}，其三维定位精度往往优于本文目标，量程也大得多。它可以
间接得到 $H$：先用反射器在基准轴上取两点拟合出轴线、再逐点巡测目标并按叉积公式
计算点到轴距离。就“能否测 $H$”而言，它是可行且精度很高的方案。真正的障碍在于
适用性：所有目标点都必须能够安放合作反射器，跟踪仪无法直接测量非合作的空间特征
点；逐点巡测需要人工把反射器依次移动到每个采样点，效率低；且设备单价高昂，难以
在现场大批量采样点的对中作业中普遍配备。因此它在实验室级少点位测量中可用，却不
适配本文面向的现场高密度对中场景。

\subsection{双目与多目立体视觉激光线重建}

学术上与本文思路最接近的是用立体视觉重建激光相关几何量的工作，也是需要重点区分
的方向。Zhong 等提出的线性三维直线重建方法先由对应点交会求线上一点、再最小化
图像角度残差求线方向，可用于双目与多目重建~\cite{zhong2021line3d}，重建出三维
直线后即可对任意点算叉积距离，与本文测量语义直接一致；但其在 30\,cm 长目标上的
距离误差为毫米量级，与本文 $15$--$50\,\mu\mathrm{m}$ 的目标相差一到两个数量级。
另有一类以正交参考物或三维定向参考物辅助双相机重建的工作，如 Xu 等借助正交参考
重建激光线投影点~\cite{xu2017ortho}、Chen 等用双相机配合三维定向参考做视觉重建
与精度分析~\cite{chen2021vision2cam}，其绝对误差约 $1\,\mathrm{mm}$。这一类方法
在“重建激光相关的三维几何量、进而可算距离”这一点上与本文可比，其共同短板是精度
停留在毫米级，无法满足工程要求。DCPAM 与之的核心区别有三：其一，采用接收屏配合
镜面反射虚像重建，把激光基准轴变为设备坐标系下可解析的三维直线，回避了从三维
空间中直接提取“看不见的激光”这一难题；其二，直接面向“任意目标点到基准轴距离”
求解，无需额外的外部参考物建立尺度；其三，精度目标比上述方法高出一到两个数量级。

\subsection{位置敏感探测器与靶板式对中仪}

工业界最主流的商用激光对中仪采用位置敏感探测器（PSD）或 CCD/CMOS 靶板：在一根轴
上安装激光发射器、在另一根轴上安装探测器，读取激光斑落在感光面上的坐标，再结合
轴的多角度旋转数据反算两轴之间的偏移与折角，代表产品如 Easy-Laser
XT770~\cite{easylaser_xt770}、Fixturlaser NXA Pro~\cite{fixturlaser_nxa} 与
Pr\"uftechnik Optalign Touch~\cite{pruftechnik_optalign}。就测量语义而言，它测的
确实是“某一点相对激光的位置”，但这一点只能是激光斑恰好击中的探测器所在位置：单点
探测分辨力虽可达微米级，却只覆盖被测的那一个点，要测另一个目标点就必须把探测器
物理搬到该点重装，无法对不在光路上、也不便安装传感器的采样点做非接触测量。因此
它虽是现役竞品，却只能逐点接触式测量、且输出的最终量是两轴的相对偏移与角度，难以
胜任对空间任意点求 $H$ 的任务。

\subsection{空白定位}

综合来看，如表~\ref{tab:compare} 所示，在“对空间任意目标点求点到激光基准轴距离”
这一测量量上，现有方案各有硬伤：语义完全匹配的钢丝找中法受下坠、振动与反复拆装
限制；精度极高的激光跟踪仪必须逐点安放合作反射器、成本高且效率低；语义直接一致的
学术双目激光线重建精度只到毫米级；现役的 PSD 靶板对中仪则只能测激光斑击中的那一
单点。真正能对空间任意点、以非接触方式、达到微米级精度地直接求 $H$ 的方法尚属
空白。DCPAM 正是要填补这一空白——用双相机加接收屏镜面虚像重建，把激光基准轴变为
设备坐标系下可解析的三维直线，从而对空间任意目标点直接解析地计算 $H$。
